% answers_short.tex —— 2005 年·普通高等学校招生考试·天津卷·文科 · 简短答案速查
% 由 tools/make_short_answers.py 自 answers.tex 的 \gksolhead/\gkanswer 提取，
% 并与 2005/2005_数学真题_天津文科_答案要点.md 逐题比对后生成，勿手改（改源 MD 后重跑）。

\begin{gknotice}
\textbf{使用说明}：本节只给最终答案便于快速核对，完整推导见解析版。答案为编写者独立推导，\textbf{非官方标准答案}。
\end{gknotice}

\gkpart{选择题}
\begin{center}\begin{tabular}{ccccccccccc}
\toprule
题号 & 1 & 2 & 3 & 4 & 5 & 6 & 7 & 8 & 9 & 10 \\
\midrule
答案 & C & A & A & A & D & C & B & A & D & B \\
\bottomrule
\end{tabular}\end{center}

\gkpart{填空题}
\begin{itemize}[leftmargin=2.6em]
  \item[\textbf{11.}] $210$。
  \item[\textbf{12.}] $2\sqrt{3}$。
  \item[\textbf{13.}] $\sqrt{2}$。
  \item[\textbf{14.}] $35$。
  \item[\textbf{15.}] $\left(-\dfrac{\sqrt{10}}{5}, \dfrac{3\sqrt{10}}{5}\right)$。
  \item[\textbf{16.}] $(-2, -1) \cup (1, 2)$。
  \item[\textbf{17.}] $\dfrac{1}{3}$。
\end{itemize}

\gkpart{解答题}
\begin{itemize}[leftmargin=2.6em]
  \item[\textbf{18.}] (1) $\sin\alpha = \dfrac{3}{5}$；(2) $\tan\left(\alpha + \dfrac{\pi}{3}\right) = \dfrac{48 - 25\sqrt{3}}{11}$
  \item[\textbf{19.}] (1) $c = 1$ 或 $c = -\dfrac{1}{2}$；(2) 当 $c = 1$ 时，$S_n = \dfrac{n(n+1)}{2}$；当 $c = -\dfrac{1}{2}$ 时，$S_n = \dfrac{4}{9} - \dfrac{6n+4}{9}\left(-\dfrac{1}{2}\right)^n$
  \item[\textbf{20.}] (1) $A_1A$ 与底面 $ABC$ 所成的角为 $60^\circ$（或 $\dfrac{\pi}{3}$）；(2) 见解析版；(3) $\dfrac{4\sqrt{3}}{27}\pi a^3$
  \item[\textbf{21.}] 此人距水平地面 $60$ 米时，视角 $\angle BPC$ 最大。
  \item[\textbf{22.}] $m$ 的取值范围为 $(-\infty, -1) \cup (4, 5] \cup [6, +\infty)$。
  \item[\textbf{23.}] (1) 焦点坐标为 $\left(0, \dfrac{1}{4a}\right)$，准线方程为 $y = -\dfrac{1}{4a}$；(2) 见解析版；(3) $y_1$ 的取值范围为 $(-\infty, -1) \cup \left(-1, -\dfrac{1}{4}\right)$
\end{itemize}
